% 整节替换 sec:numerical-method，结束位置为 sec:numerical-results 之前。
% 依赖原有 NumTheoryRef、NumIntervals、NumHorizon、NumDegree、NumMeshDetails、
% NumFindingParagraph、NumCheckStatement，以及本计划新增的 NumRobustnessStatement。
% 没有实验时，结果宏为空；不把下面的方法说明变成“检查已通过”的结论。
\section{数值方法与实验设置}
\label{sec:numerical-method}

本节围绕不同初值区域，比较解析最优策略与直接配点输出，
重点考察等待阶段、容量弧和完全跟踪阶段的出现顺序及其连接位置。
算例组织借鉴 \citet[第6节]{AvramFreddiGoreac2022} 的区域相关情景设计，
但使用本文的状态方程、控制成本及
\NumTheoryRef{def:regions}{分区定义}。
该文献的有限时域问题包含目标约束，
本文仅借鉴其几何展示方式，不将其结果作为下述无终端约束截断的等价性依据。
全局最优性与几乎处处唯一性分别由
\NumTheoryRef{thm:verification}{验证定理}和
\NumTheoryRef{thm:unique-q}{唯一性定理}给出；
数值实验用于检验这些解析结论的计算一致性。

五个主算例采用相同参数
\begin{equation}
 p=0.5,\qquad c=2,\qquad\gamma=0.3,\qquad K=0.15,
 \label{eq:numerical-common-parameters}
\end{equation}
仅改变初始状态。
$s,i$ 为相对于初始总量的比例，允许 $s_0+i_0<1$；
此时 $h=0.3$、$\ell=r=0.15$。
这些参数用于展示理论结构，不作为特定疫情的参数标定。

直接配点处理如下固定时域问题：
\begin{equation}
\begin{aligned}
 \min_q\quad &J_T(q)=\int_0^T pcq(t)\,\mathrm dt,\\
 \text{满足}\quad
 &\dot s=-c[p+(1-p)q]si,\qquad
 \dot i=[pc(1-q)s-\gamma]i,\\
 &0\le q\le1,\quad 0\le s\le1,\quad0\le i\le K,\\
 &(s(0),i(0))=(s_0,i_0).
\end{aligned}
\label{eq:numerical-ocp}
\end{equation}
该问题不增加终端状态约束、终端成本或二次控制正则化，
也不固定阶段数、切换时刻或容量控制。
解析策略可用于初始化，但不进入数值优化的目标或额外约束。

计算方法为 MATLAB/OpenOCL 直接配点\citep{Koenemann2019OpenOCL}，
通过 CasADi\citep{Andersson2019CasADi} 构造离散问题，
由 IPOPT 求解非线性规划\citep{WachterBiegler2006}。
基准设置为 $T=\NumHorizon$、$N=\NumIntervals$、$d=\NumDegree$；
控制在每个网格区间内为常数。
\NumMeshDetails
控制成本在完整求解时域上按实际区间长度计算：
\begin{equation}
 J_{\mathrm{OCL}}=pc\sum_{k=1}^{N}q_k\Delta t_k.
 \label{eq:numerical-cost-sum}
\end{equation}
绘图窗口不作为成本积分范围，求解输出不通过平滑或裁剪替换。

对于每个待报告结果，用原始分段常数控制重新积分状态方程，
检查区间内部的容量峰值、与配点状态的偏差以及终端零控制延拓。
对正状态，未来无控制峰值为
\[
 P_0(s,i)=
 \begin{cases}
  i,&s\le h,\\
  i+s-h-h\ln(s/h),&s>h.
 \end{cases}
\]
终端延拓检查对重积分得到的状态使用 $P_0(s(T),i(T))$，
同时检查终点前一段时间内的控制是否已接近零。
这些是求解后的浮点诊断，不进入问题\eqref{eq:numerical-ocp}。
它们不单独证明连续时间严格可行性，
也不将有限时域离散问题与原无穷时域问题视为严格等价。

阶段和过渡区间首先依据数值控制及容量接近程度识别，
再与解析事件比较。
跳跃附近按过渡网格区间报告位置，
不以首个纯控制单元的起点代替连续切换时刻。
附加复核分别改变初始猜测、网格和终止时刻；
网格比较固定 $T$，时域比较固定控制区间宽度，
以区分离散误差和截断效应。
详细复核保留在配套结果文件中，正文仅保留与阶段结构判断相关的摘要。
局部求解器的重复收敛及有限样本复核均不替代解析证明。

% 下列宏由实际保存数据生成；未完成时保持为空。
\NumFindingParagraph{\NumCheckStatement}
\NumFindingParagraph{\NumRobustnessStatement}
